\documentclass{article}
\usepackage[brazil]{babel}

\usepackage[margin=1in]{geometry}
\usepackage{amsmath}
\usepackage{tikz}

\title{Processos Estocásticos}
\author{Thiago Tomás de Paula}
\date{4 de maio de 2026}

\begin{document}

\maketitle

\section*{\sffamily Questão 1}
Um ``Coffe Shop'' tem 4 tipos diferentes de café. Você pode pedir seu café em uma xícara
pequena, média ou grande. Você também pode optar por adicionar creme, açúcar ou leite (qualquer
combinação é possível. Por exemplo, você pode optar por adicionar os três). 

\begin{itemize}
    \item [a)] 
    De quantas maneiras você pode pedir seu café?
    
    \item [b)] 
    Durante quatro vezes consecutivas você visita esse “Coffe Shop” com um amigo e pede para ele escolher aleatoriamente um tipo de café. Qual é a probabilidade dele não escolher, pelo menos uma vez: em uma xícara pequena e adicionar simultaneamente creme, açúcar e leite

    \item [c)]
    Para a mesma situação da letra (b), qual é a probabilidade dele não escolher pelo menos uma vez: xícara média, café com leite e açúcar?
\end{itemize}
\rule{\linewidth}{0.5pt}

\noindent\textbf{\sffamily Respostas}
\begin{itemize}
    \item [a)] 
    Para os adicionais, existem $3\times 2\times 1 = 3! = 6$ combinações possíveis, assumindo que uma opção não pode ser repetida. 
    Temos então
    %
    \begin{equation*}
    \boxed{
        \underbrace{4}_{\text{Tipos de café}} \times 
        \underbrace{3}_{\text{Tamanhos de xícara}} \times 
        \underbrace{6}_{\text{Adicionais}} 
        = 72
    }
    \end{equation*}
    %
    maneiras de se pedir o café.
    
    \item [b)] 
    Existem apenas 4 maneiras de se pedir o café impondo que a xícara seja pequena \emph{e} que todos os adicionais estejam presentes, uma para cada tipo de café. 
    A probabilidade de tal pedido ser feito é então 
    %
    \begin{equation*}
        p_\neg = 
        \frac{4}{72} =
        \frac{1}{18}
    \end{equation*}
    %
    de forma que a chance desse pedido \emph{nunca} ser feito mesmo em \emph{quatro} vezes consecutivas se torna
    %
    %
    \begin{equation*}
    \boxed{
        p = 
        (1 - p_\neg)^4 = 
        \left(\frac{17}{18}\right)^4 \approx
        79,56\%
    }
    \end{equation*}
    

    \item [c)]
    Para a mesma situação da letra (b), qual é a probabilidade dele não escolher pelo menos uma vez: xícara média, café com leite e açúcar?

    O pedido com xícara média, leite e açúcar pode ser feito de 8 maneiras diferentes, isto é, adicionando ou não o creme para cada um dos 4 tipos de café. A probabilidade desse pedido ser feito é então
    %
    \begin{equation*}
        p_\neg = 
        \frac{8}{72} =
        \frac{1}{9}
    \end{equation*}
    %
    de forma que a chance desse pedido \emph{nunca} ser feito mesmo em \emph{quatro} vezes consecutivas se torna
    %
    %
    \begin{equation*}
    \boxed{
        p = 
        (1 - p_\neg)^4 = 
        \left(\frac{8}{9}\right)^4 \approx
        62,43\%
    }
    \end{equation*}
\end{itemize}

\newpage

\section*{\sffamily Questão 2}
Considere um experimento contendo quatro cartas marcadas com os números inteiros 1 a 4. No
experimento se sorteia aleatoriamente um par de cartas da sacola, ou seja, cada par de carta constitui uma saída do experimento randômico.

\begin{itemize}
\item [a)] Encontre o espaço de amostra $S_1$, para o experimento, se a primeira carta for retornada para a sacola antes que a segunda seja sorteada.

\item [b)] Em $S_1$, qual é a probabilidade de serem sorteadas as cartas com os números (1, 2) ou (2,3) ou (3,4) em uma saída do experimento aleatório.

\item [c)] Em $S_1$ qual é a probabilidade de ser sorteada a carta número 3 dado que a carta1 já foi sorteada?

\item [d)] Em $S_1$, qual é a probabilidade de não serem sorteadas sucessivamente as cartas com os números (2, 4)
e (1,3).

\item [e)] Encontre o espaço de amostral $S_2$, para experimento, se o primeiro cartão não for retornado para a
sacola.

\item [f)] Para o espaço amostral $S_2$, qual a probabilidade da carta de ser sorteado o par (3, 3)?
\item [g)] Em $S_2$, qual a probabilidade de não serem sorteadas os pares (1,2) e (2,1)?
\item [h)] Em $S_2$, qual probabilidade da carta de número 4 não ser sorteada?
\end{itemize}
\rule{\linewidth}{0.5pt}

\textbf{\sffamily Respostas}

\begin{itemize}

\item [a)] Havendo reposição, tanto o sorteio da primeira carta quanto o da segunda têm 4 resultados. 
Sendo os saques eventos independentes, obtemos $S_1$ pela multiplicação dessas quantidades:
%
\begin{equation*}
\boxed{
    S_1 = 4 \times 4 = 16
}
\end{equation*}

\item [b)] A probabilidade de se observar qualquer par em particular é 
%
\begin{equation*}
\boxed{
\frac{1}{S_1} = \frac{1}{16} 
              \approx 6.25\%
}
\end{equation*}

\item [c)] Como a carta 3 é uma das 4 opções equiprováveis, a probabilidade condicional é de 
$1/4 = 25\%$. 

\item [d)] Sendo os resultados consecutivos eventos independentes e a probabilidade de cada resultado $1/16$, obtemos a resposta pelo produto dessa probabilidades:
%
\begin{equation*}
\boxed{
    p = \frac{1}{S_1} \times \frac{1}{S_1}
      = \frac{1}{256}
      \approx 0.39\%
}
\end{equation*}

\item [e)] Nesse caso, existem 4 opções para o primeiro cartão retirado, e 3 para o segundo. Tomando o produto, conclui-se que
%
\begin{equation*}
\boxed{
    S_2 = 4 \times 3 
        = 12
}
\end{equation*}

\item [f)] É impossível sortear um par com números iguais, então a probabilidade é 0.
\item [g)] A probabilidade de um par em particular ser sorteado é $p_\neg = 1/S_2 = 1/12$. 
Dessa forma, a probabilidade dos pares $(1,2)$ e $(2,1)$ \emph{não} serem sorteados é 
%
\begin{equation*}
\boxed{
    p = 1 - p_\neg
      = \frac{11}{12}
      \approx 91,67\%
}
\end{equation*}

\item [h)] A probabilidade de um número ser sorteado é 
$p_{\neg, 1} = 1/4 \approx 25\%$
no primeiro saque e 
$p_{\neq, 2} = \frac{1}{3}$ no segundo.
Logo, a probabilidade de $4$ não ser pego nem no primeiro saque nem no segundo é
%
\begin{equation*}
\boxed{
    p  = (1 - p_{\neg, 1}) (1 - p_{\neg, 2})
       = \frac{3}{4} \times \frac{2}{3}
       = \frac{1}{2}
       = 50\%
}
\end{equation*}
\end{itemize}

\section*{\sffamily Questão 3}
Um experimento consiste em rolar um dado até obter um 6 com o máximo de 3 jogadas.

\begin{itemize}
\item [a)] Determine o espaço amostral $S_1$ onde estão presentes todas as possibilidades de interesse.

\item [b)] Qual a probabilidade de após três jogadas do dado, nenhum elemento do espaço amostral $S_1$ acontecer?
\end{itemize}
\rule{\linewidth}{0.5pt}

\textbf{\sffamily Respostas}

\begin{itemize}
\item [a)] 
Assumindo que o dado deve ser rolado mesmo se um 6 for obtido, existem
\begin{itemize}
    \item $6 \times 6 = 36$ maneiras de se obter 6 no primeiro lance, sendo o 2° e 3° quaisquer valores;
    \item $5 \times 6 = 30$ maneiras de se obter 6 no segundo lance, sendo o 1° lance diferente de 6;
    \item $5\times 5 = 25$ maneiras de obter 6 apenas no 3° lance.
\end{itemize}
Somando as maneiras independentes, temos
%
\begin{equation*}
\boxed{
    S_1 = 36 + 30 + 25 
        = 91
}
\end{equation*}

\item [b)] 
Existem $6 \times 6 \times 6 = 216$ resultados possíveis para o caso em que o dado é rolado exatamente 3 vezes, dos quais $S_2=91$ constituem um jogo vitorioso.
A probabilidade de perder o jogo é então
%
\begin{equation*}
\boxed{
    p = \frac{216 - S_1}{216}
      = \frac{125}{216}
      \approx 57,87\%
}
\end{equation*}
\end{itemize}

\section*{\sffamily Questão 4}
Um experimento aleatório consiste em jogar 3 dados distintos. Contudo, um dos dados tem 1 (um)
único ponto desenhado em todas as suas faces.

\begin{itemize}
\item [a)] Determine o espaço amostral do experimento.
\item [b)] Qual é a probabilidade de um evento A no qual a somas dos três dados é igual a 8?
\item [c)] Qual é a probabilidade de um evento B no qual a somas dos três dados é superior a 10?
\item [d)] Encontre a probabilidade de um evento C onde a soma dos três dados é superior a 12.
\end{itemize}
\rule{\linewidth}{0.5pt}

\textbf{\sffamily Respostas}

\begin{itemize}
\item [a)] 
Os pares de valores dos dados comuns se combinam em $6\times 6 = 36$ maneiras distintas. 
Essas combinações ocorrem para cada uma das 3 posições que o dado incomum ocupa (primeiro, segundo e terceiro a ser lançado), de forma que o conjunto amostral tem tamanho
%
\begin{equation*}
\boxed{
    S = 3\times 36 
      = 108 
}
\end{equation*}


\item [b)] 
Em $A$, os dados comuns retornam os pares 
$(6, 1)$,
$(5, 2)$,
$(4, 3)$,
$(3, 4)$,
$(2, 5)$ ou
$(1, 6)$.
Cada um destes 6 pares podem ocorrer para cada uma das 3 posições do dado fixo, de forma que existem $6\times 3 = 18$ maneiras de se obter lances com soma 8.
Logo,
%
\begin{equation*}
\boxed{
    p(A) = \frac{18}{108}
         \approx 16,67\%
}
\end{equation*}

\item [c)]
Em $B$, os dados comuns retornam os pares 
$(6, 4)$,
$(6, 5)$,
$(6, 6)$,
$(4, 6)$ ou
$(5, 6)$.
Cada um destes 5 pares podem ocorrer para cada uma das 3 posições do dado fixo, de forma que existem $3\times 5 = 15$ maneiras de se obter lances com soma 11 ou maior.
Logo,
%
\begin{equation*}
\boxed{
    p(B) = \frac{15}{108}
         \approx 13,88\%
}
\end{equation*}

\item [d)]
Em $C$, os dados comuns retornam o par 
$(6, 6)$.
Esse par pode ocorrer para cada uma das 3 posições do dado fixo, de forma que existem $3\times 1 = 3$ maneiras de se obter lances com soma 13.
Logo,
%
\begin{equation*}
\boxed{
    p(C) = \frac{3}{108}
         \approx 2,78\%
}
\end{equation*}
\end{itemize}

\section*{\sffamily Questão 5}
Uma concessionária revendedora de automóveis oferece veículos com as seguintes opções: 

\begin{itemize}
\item [1)] Com ou sem transmissão automática; 2) Com ou sem ar-condicionado; 
\item [2)] Com ou sem GPS; 
\item [3)] Com duas possíveis opções de um sistema de reprodução áudio estéreo; 
\item [4)] Com quatro cores padrão.
\end{itemize}

\begin{itemize}
\item [a)]
Se o espaço de amostra consistir no conjunto de todos os tipos de veículos possíveis, qual é o número
de resultados possíveis do espaço amostral?

\item [b)]
Desenhe o diagrama em árvore para os diferentes tipos de veículos.

\item [c)]
Se cada tipo de automóvel tem a mesma probabilidade de ser comercializado pela concessionária,
então, determine a probabilidade do automóvel vendido ter GPS.

\item [d)]
Qual a probabilidade do automóvel vendido ter GPS dado que ele possui transmissão automática?
\end{itemize}
\rule{\linewidth}{0.5pt}

\textbf{\sffamily Respostas}

\begin{itemize}
\item [a)]
Multiplicando as maneiras de se configurar os elementos em separado, obtemos
%
\begin{equation*}
\boxed{
    \underbrace{2}_{\text{Transmissões}} \times
    \underbrace{2}_{\text{AC}} \times
    \underbrace{2}_{\text{GPS}} \times
    \underbrace{2}_{\text{Áudio}} \times
    \underbrace{4}_{\text{Cores}}
    =
    64
}
\end{equation*}

\item [b)]
\begin{center}
    \begin{tikzpicture}[> = latex]
        \foreach \t in {45, -45}{
            \draw (0, 0) -- (\t:2.5) node (prev0) {};
            
            \foreach \ac in {30, -30}{
                \draw (prev0) --++ (\ac:1.75) node (prev1) {};
                
                \foreach \gps in {15, -15}{
                    \draw (prev1) --++ (\gps:1.5) node (prev2) {};
                
                    \foreach \a in {10, -10}{
                        \draw (prev2) --++ (\a:1.25) node (prev3) {};
                        
                        \foreach \c in {10, 3, -3, -10}{
                            \draw (prev3) --++ (\c:1) node (prev4) {};
                        }
                    }
                }
            }
        }
        \draw [->] (prev0) |-++ (-0.5, -0.5) node [anchor=east] {Transmissão};
        \draw [->] (prev1) |-++ (-0.5, -0.5) node [anchor=east] {AC};
        \draw [->] (prev2) |-++ (-0.5, -0.5) node [anchor=east] {GPS};
        \draw [->] (prev3) |-++ (-0.5, -0.5) node [anchor=east] {Áudio};
        \draw [->] (prev4) --++ (0, -0.5) node [anchor=north] {Cores};
    \end{tikzpicture}
\end{center}

\item [c)]
Como a presença e ausência de GPS são igualmente prováveis, em média metade dos carros medidos terão GPS, ou 50\%.

\item [d)]
Como as opções são independentes, a probabilidade continua de 50\% como no item (c).
\end{itemize}


\end{document}

